\documentclass[final,a4paper,12pt]{article}

\usepackage{notes}
\setlength\parskip{1ex}
\setlength\parindent{0cm}
\usepackage{float}
\usepackage{mathtools}
\usepackage{cancel}
\usepackage{amsthm}
\theoremstyle{definition}
\newtheorem{example}{Example}[section]
\newtheorem{problem}{Problem}[section]

\newcommand\qeq{\stackrel{\mathclap{\normalfont\mbox{?}}}{=}}

\title{Geophysics III}
\author{Algebra Tutorial}
\date{}

\begin{document}
\maketitle

\section{Introduction}
Algebra is a set of rules and theorems that govern a broad swath of mathematics.
While we'll employ calculus as well, much of what we do will be
substitution and rearangement of abstract algebraic expressions.  Abstract
algebra simply means the use of symbols to represent numbers and concepts rather
than numbers themselves.  We'll use this absract representation so that we can
learn about the general patterns rather than the solution for any specific case.
And therein lies the power of abstract algebra: the ability to write a general
expression applicable to many, perhaps infinitely many, cases.

This tutorial is meant to be a reminder of basic algebraic functions,
properties, and techniques.  If you have never had math beyond algebra, it may
introduce a few new ideas and concepts that you weren't exposed, but are well
within the confines of algebra.  To some this may seem trivial, but to others it
will be a useful review.


\section{Rearranging equations}
\subsection{Rearranging rules}
There are a few simple rules that you must remember when rearranging equations.
You'll use them frequently when performing derivations.  The rules:
\begin{enumerate}
    \item perform the same operation to both sides of the equation;
    \item when adding or subtracting, make sure the sum is zero since adding
    zero doesn't change a result; and
    \item when multiplying or dividing, make sure the product is one because
    multiplying anything leaves the result unchanged;
    \item when raising an exponent or taking a root, make sure the power is
    one.
\end{enumerate}

{\it Practice:} \\
Show $2 r (x^2 - 1) = x^2 - 2$ can be written as
\begin{equation}
    x^2 = \frac{2(1 - r)}{(1 - 2r)} .
\end{equation}

\subsection{Common mistakes}
Most of the common errors occur when dealing with fractions.  Although the rules
of dealing with fractions are quite simple, the errors seem to happen because of
temptation; temptation to cancel terms that simplify the expression, but
generally too much.  If you are ever unsure about a rule, TEST IT!  Put real
numbers in and test the before and after (use simple numbers).  If the before
and after don't give the same result, then what you did is invalid.  If it does,
try another set and if it still works then you probably did it right.  I'll give
a few examples of things that I commonly see but are incorrect.

\begin{example}
    Mistakes commonly occur when trying to cancel terms that show up in both the
    numerator and denominator.  For instance,
    \begin{equation*}
        \frac{\cancel{z^2}}{x^2 + \cancel{z^2}} \neq \frac{1}{x^2} ,
    \end{equation*}
    where the $z^2$ terms cannot be canceled.  If we put in some numbers, 1 for
    $x$ and 2 for $z$,
    \begin{align*}
        \frac{2^2}{1^2 + 2^2} &\neq \frac{1}{1^2} \\
        \frac{4}{5} &\neq 1 .\\
    \end{align*}
    It turns out that the above expression cannot be simplified further
    
    We cannot this way either
    \begin{equation*}
        \frac{x^2 + \cancel{z^2}}{\cancel{z^2}} \neq x^2 ,
    \end{equation*}
    which we can test as above,
    \begin{align*}
        \frac{1^2 + 2^2}{2^2} &\neq 1^2 \\
        \frac{5}{4} &\neq 1 .
    \end{align*}

    You can partially cancel the above expression,
    \begin{align*}
        \frac{x^2 + z^2}{z^2} &= \frac{x^2}{z^2} + \cancelto{1}{\frac{z^2}{z^2}} \\
        &= \frac{x^2}{z^2} + 1 \\
    \end{align*}
    Using our inputs above you can verify that the result is $1+\frac{1}{4} =
    \frac{5}{4}$.
\end{example}

\begin{example}
    \begin{equation*}
        \frac{x^2 - 2 z^2}{x^2} - \frac{x^2}{z^2} \neq \frac{-2 z^2}{x^2}
    \end{equation*}
\end{example}

\clearpage

\subsection{Rules for special functions} 
There are a number of handy properties of logs and exponents that seem to be
commonly forgotten, which I don't blame you for.  It is quite likely that you
haven't used them since you learned them the first time.  Turns out, we use them
in science quite frequently; although, if you have only ever been given
equations you my not realize that it has been done.  These rules can be useful
for simplifying equations or reorganizing them when performing derivations.

\subsubsection{Logarithms}
We can write a logarithm generally as
\begin{equation*}
    \log_{
        \mlbox[labelstyle={below of=m\themathLableNode,below of=m\themathLableNode}]
        {a}{base}}
        \mlbox[]{u}{argument} .
\end{equation*}
Although the logarithm above is written generically, the bases used is in most
often in physics are either 10 or $e$, the natural number.  For $\log_10$, we
often use a shorthand $\Log$ and sometimes $\log$.  For $\log_e$, the shorthand
$\ln$ or $\log$ are often used.  Because $\log$ is used as the shorthand for
both bases 10 and $e$ in different texts, you will need to identify which base
$\log$ applies.

Properties of logarithms, $u, v > 0$,
\begin{gather*}
    \log_a a = 1 \\
    \log_a a^n = n \\
    \log_a u^n = n \log_a u \\
    \log_a u v = \log_a u + \log_a v \\
    \log_a \frac{u}{v} = \log_a u - \log_a v
\end{gather*}
To change base, i.e., from $a$ to $b$ (where $a, b \neq 1$),
\begin{equation*}
    \log_b u = \frac{\log_a u}{\log_a b}.
\end{equation*}
There are several key points on a logarithm that are useful to know,
\begin{align*}
    \Log 0 &= -\infty & \log 0 &= -\infty \\
    \Log 1 &= 0 & \log 1 &= 0 \\
    \Log 10 &= 1 & \log e &= 1 .
\end{align*}


\subsection{Exponentials}
Hidden in the relationships above is the relationship between logarithms and
exponentials as the inverse operations of each other, i.e.,
\begin{align*}
    \log_a a^n &= n \\
    a^{\log_a n} & = n .
\end{align*}
Properties of exponentials,
\begin{align*}
    a^m a^n &= a^{m + n} \\
    \frac{a^m}{a^n} &= a^{m - n} \\
    (a^m)^n &= a^{mn} \\
    a^{-n} = \frac{1}{a^{n}} \\
    a^0 = 0 \\
    a^1 = a .
\end{align*}
If $a$ and $n$ are positive, $a^{\frac{1}{n}} = \sqrt[n]{a}$ .

One must be carful with negative powers when $a$ is also negative as they may
yield a complex rather than real number (see complex number tutorial).


\subsection{Simplifying solutions}
We will try to write our solutions in as simple and elegant a form as possible.
One reason to do so is that it can be easier to interpret, compute, and/or graph.
But what one decides as an elegant representation is subject to debate.  Also,
it may be easier for a computer to compute a less elegant representation.
Nevertheless, we will strive to do so.  Here are a few suggestions to derive
elegant solutions:
\begin{itemize}
    \item combine like terms, no use carrying around the
    extra weight, e.g., $3x + x = 4x$ ;

    \item reduce fractions, again
    reduces extra weight, e.g., $\frac{6x^2}{64x} = \frac{3x}{32}$;

    \item find least common multiples, this can help us simplify a solution,
    e.g., $\frac{1}{x_1} + \frac{1}{x_2} = \frac{x_2 + x_1}{x_1 x_2}$ ;

    \item factorize polynomials,, doing so can help us quickly identify roots
    (zero crossings),  e.g., $\rho_1 r - \rho_2 r = (\rho_1 - \rho_2) r$ ; and

    \item expand/distribute (opposite of factoring), sometimes this simplifies
    the expression, e.g., $3z^2 - (x^2 - z^2) = 4z^2 - x^2$ ; and

    \item write linear relationships in slope-intercept form---the usefulness
    speaks for itself---or non-linear equations in a standard form, e.g.,
    Pythagorean equation.
\end{itemize}
Also, for equations containing complex numbers, it is best to write in a form
without complex numbers in the denominator, e.g., $\frac{1}{1 - i} = \frac{1}{1
- i}\frac{1 + i}{1 + i} = \frac{1 + i}{2}$.

Since some of the above suggestions for simplifying equations are contradictory,
{\it how will you know what I'm looking for as a grader?}  The key is having all
the correct elements so there is a fair amount of latitude to your answer so
long as it has been simplified.


\subsection{Practice problems}
\begin{problem}
    If
    \begin{equation*}
        \sin \theta = \frac{x}{h},
    \end{equation*}
    show
    \begin{equation*}
        \cos \theta = 1 - \frac{x^2}{h^2}.
    \end{equation*}
    You'll need to use the Pythagorean Theorem with respect to the sides of a
    right triangle with one of the two remaining angles $\theta$.
\end{problem}

\begin{problem}
    Show that
    \begin{equation*}
        y = 2 \left( \frac{1}{x_1^2} - \frac{1}{x_2^2} \right)^{1/2}
    \end{equation*}
    is equivalent to
    \begin{equation*}
        y = \frac{2 (x_2^2 - x_1^2)^{1/2}}{x_1 x_2} .
    \end{equation*}
\end{problem}

\begin{problem}
    Show that
    \begin{equation*}
        b = a_1^{v_1} a_2^{v_2}
    \end{equation*}
    is equivalent to
    \begin{equation*}
        \log b = v_1 \log a_1 + v_2 \log a_2 .
    \end{equation*}
\end{problem}


\section{Special operators}
\subsection{Summation, $\sum$}
\begin{equation*}
    \sum_{
    \mlbox[
        labelstyle={below of=m\themathLableNode,below of=m\themathLableNode,xshift=-1.5cm},
        linestyle={out=90,in=270, latex-}]
        {i}{index}
    =\mlbox[labelstyle={below of=m\themathLableNode,below of=m\themathLableNode}]{a}{lower limit}
    }^{\mlbox[]{N}{upper limit}}
    \mlbox[
        labelstyle={right of=m\themathLableNode,right of=m\themathLableNode}]
        {x_i}{argument}
\end{equation*}
At times we will want to add or multiply a known, and regular set of discrete
numbers (note this will eventually lead to calculus).  These will show up from
time to time throughout the course.  For example, suppose we
wanted to add the integers from 1 to 5 whose result is 15,
\begin{equation}
    1 + 2 + 3 + 4 + 5 = 15.
\end{equation}
Because mathematicians are lazy, they seek a shorthand to write this in a
simpler way.  Rather than add up all the numbers at once, let's group these
numbers to perform only one operation at a time, i.e.,
\begin{equation}
    (((\,(1) + 2) + 3) + 4) + 5.
\end{equation}
Working with the inner most parentheses first, we find 1; inside the next set we
find $(1 + 2)$, which yields 3.  Continuing each step, we finally arrive at
15.  This is the point where the simplifying notation comes in.  Because each of
the above operations is an addition, we use the capital Greek letter sigma
($\Sigma$) to denote the sum.  We also give the sum limits: where to start
(above) and where to end (below).  Thus, we can write the sum as
\begin{equation}
    \sum_{i = 1}^5 i \quad \text{where } i = 1,2,\ldots,5,
\end{equation}
In this case, we perform the summation over the values of the index, $i$, from 1
to 5, incrementing $i$ by a single value each step.

The increments of the index, need not be unity, but could be anything we need
them to be so long it is regular or defined in some ordered way.  The argument,
also $i$ in the example above, could also be anything.

We can express the arithemetic mean using the summation.  For a set of $N$
numbers, the arithmetic mean is given as
\begin{equation}
    \bar{x} = \frac{x_1 + x_2 + \ldots + x_N}{N},
\end{equation}
or in short form,
\begin{equation}
    \bar{x} = \frac{1}{N}\sum_{i=1}^{N} x_i .
\end{equation}
Likewise, the sample standard deviation is defined as
\begin{equation}
    \sigma_{x} = \left[
        \frac{1}{N-1} \sum_{i=1}^{N}(x_{i} - \bar{x})^2
        \right]^{1/2}.
\end{equation}

{\it Practice:} \\
Prove to yourself
\begin{equation}
    \sum_{i = 1}^4 i^2 = 30
\end{equation}

\subsection{Product, $\prod$}
There is an analagous operator for multiplication, just as we used the $\Sigma$
for summation, the capital Greek letter pi ($\Pi$) is used for the product of a
sequence.  For example, we can write the product of $N$ numbers as
\begin{equation}
    \prod_{i = 1}^N x_i = x_1 \cdot x_2 \cdot x_3, \ldots, x_N .
\end{equation}

This notation arises when we use a geometric mean
\begin{equation}
    \bar{x}_g = \left( \prod_{i = 1}^N x_i \right)^{1/N} .
\end{equation}

{\it Practice:} \\
Prove to yourself
\begin{equation}
    \prod_{i = 1}^4 i^2 = 576.
\end{equation}


\section{Common functions}
There are a number of functions that are quite common
(Figure~\ref{fig:functions}).  We will use nearly all of these functions at
least once if not repeatedly within the course.  It is extremely important to
know what these functions generally look like and know their attributes.  The
attributes that are the most important to understand regard how they can be
manipulated to fit physical models and/or data and provide insight into the
physical processes that we use them to describe.  These attributes include their
scale, amplitude, slope, and key points: location of asymptotes, zero crossings,
intercepts etc.
\begin{figure}[tbhp]
  \centerline{
    \includegraphics[angle=90]{common_functions.eps}
  }
  \caption{Common functions.}
  \label{fig:functions}
\end{figure}

\subsection{Linear}
A line is one of the simplest functions (and most useful) functions that we can
use in any science.  A linear process is one in which one variable is
proportional to another by a constant.  This constant provides the
scale---commonly know as the slope.  Add an additional constant that accounts
for a shift or static offset between each of the variables (the intercept).  The
equation of a line is given by
\begin{equation}
    y = m x + b ,
\end{equation}
where $m$ is the slope and $b$ is the intercept.  Depending on our process the
slope and/or intercept may be represented by multiple physical constants that
combine to make a single constant.  An alternative form of the linear equation
that is at times handy can be written,
\begin{equation}
    y - y_0 = m (x - x_0).
\end{equation}
A line also represents a first order, and the simplest, polynomial.  You will
see lines in many places throughout geophysics.

\subsection{Parabola/Quadratic}
A parabola is a second order polynomial, written
\begin{equation}
    y = a_2 x^2 + a_1 x + a_0 ,
\end{equation}
with three constants, $a_0$, $a_1$, and $a_2$.  You'll notice that the form I
have written in Figure~\ref{fig:functions} as
\begin{equation}
    y = (x - r_1)(x - r_2) ,
\end{equation}
where $r_1$ and $r_2$ are the roots (zero crossings) of the parabola.  While the
equation can always be written in this form, $r_1$ and $r_2$ are only real
valued when the equations intersect the $x$-axis.  When the parabola does not
intersect the $x$-axis, the roots are complex numbers (we'll get to complex
numbers later).  To find the roots of a parabola, we can use the famous
quadratic formula,
\begin{equation}
    r = \frac{-b \pm \sqrt{b^2 - 4ac}}{2a} .
\end{equation}
You'll see a quadratic when we get to heat flow.

\subsection{Hyperbola}
A hyperbola is a conic section (as are parabolas, circles, and ellipses).  They
are also the shape of a seismic reflection.  A hyperbola can be expressed in the
form
\begin{equation}
    y = \pm a \left[1 + \frac{x^2}{b^2} \right]^{1/2} ,
\end{equation}
where $a$ and $b$ are constants.  A hyperbola has two asymptotes, each with a
slope defined by the ratio of the constants, i.e., $m = \pm\frac{a}{b}$.  The
foci are also defined by the constants, $c^2 = a^2 + b^2$.  The $y$-axis
intercept of the ellipse is defined by the constant, $a$.

\subsection{Inverse Power Laws}
Functions which are represented by the inverse of a power,
\begin{equation}
    y = x^{-n} ,
\end{equation}
where $n$ is a positive multiple of $1/2$ are relatively common.  For instance,
gravitational acceleration is proportional to the inverse square of the
distance between two bodies, i.e., $\propto r^{-2}$.  As $x \to 0$,
the function grows to infinite magnitude, and as $x \to \infty$, the function
approaches zero.

Now it is your turn. For the remainder of these, I'll let you have some practice
writing up your own observations.  I suggest you graph these functions.  Change
some of the constants and see how it changes the graph.  What attribute(s) did
it affect?

{\it Remember, these functions are used to describe physical processes and many of
the key properties can be determined directly from the function and/or graph.}

\subsection{Exponentals}
Where might these be common?  What is the general form?  Are there any key
points, asymptotes, discontinuities?
\vspace{16ex}

\subsection{Logarithm}
\vspace{16ex}

\subsection{Gaussian}
\vspace{16ex}

\subsection{Error function}
\vspace{16ex}

\subsection{Trigonometric functions}
\subsubsection{Sine}
\vspace{16ex}
\subsubsection{Cosine}
\vspace{16ex}
\subsubsection{Tangent}
\vspace{16ex}

\subsection{Delta function}
A delta function is expresed as
\begin{equation}
    y = \delta (x - x_0).
\end{equation}
The delta function is a special function that is always zero except when $x =
x_0$ at which point the function is unity.  While a bit strange, these are very
useful functions when we work with waves, but they are used in many engineering
and physical fields.

\subsection{Heaviside function}
\vspace{16ex}

\end{document}
